\section{Related Works}

\paragraph{Explicit reasoning in token space.}
Chain-of-thought prompting demonstrated that encouraging LLMs to generate intermediate reasoning steps substantially improves performance on complex tasks \cite{wei2022chain}. Subsequent work proposed refinements such as decomposed prompting \cite{khot2022decomposed}, least-to-most prompting \cite{zhou2022least}, and self-consistency or guided beam search \cite{xie2023self} to improve robustness. Tree-based and search-oriented approaches explicitly explore multiple reasoning paths during decoding \cite{yao2023tree}. Despite their success, these methods remain fundamentally constrained by autoregressive token generation and must commit to discrete reasoning paths sequentially. Related work also frames LM reasoning as search and planning, either by learning search dynamics or studying planning failures at scale \cite{gandhi2024stream,lehnert2024beyond,hao2023reasoning}.

\paragraph{Efficiency-oriented CoT variants.}
A parallel line of work focuses on reducing the cost of explicit reasoning without changing its discrete nature. Methods such as concise or draft-based reasoning \cite{xu2025chain,madaan2022text}, token skipping and early exit strategies \cite{xia2025tokenskip,yang2025dynamic}, and sketch-based reasoning \cite{aytes2025sketch} aim to write fewer tokens or terminate reasoning early. Surveys on efficient reasoning systems highlight that these approaches primarily optimize surface-level generation rather than internal computation \cite{feng2025efficient,sui2025stop}.

\paragraph{Theoretical analyses of chain-of-thought.}
Several works study why CoT improves reasoning from a representational or expressivity perspective. Prior analyses show that CoT can increase the effective depth and expressive power of Transformers on serial problems \cite{merrill2023expressive,li2024chain,feng2023towards}, but also highlight limitations arising from greedy token-level commitment \cite{saparov2022language}. These results motivate moving beyond discrete reasoning substrates.

\paragraph{Implicit and latent reasoning with tokens.}
Early attempts to internalize reasoning include pause tokens and filler tokens that allow models to allocate additional computation without explicit output \cite{goyal2023think}, as well as methods that distill explicit CoT into implicit representations \cite{deng2024explicit,deng2024implicit}. Other approaches probe or guide hidden computation without fully abandoning discrete token interfaces \cite{pfau2024lets,wang2024guiding}. While these methods blur the boundary between language and computation, they retain token-level structure and do not fully exploit continuous latent spaces.

\paragraph{Continuous latent reasoning.}
Coconut introduced continuous latent thoughts as a replacement for explicit rationale tokens, enabling models to reason in hidden state space while maintaining standard autoregressive decoding \cite{hao2024training}. Follow-up work analyzed the emergence and training dynamics of latent superposition states \cite{zhu2025emergence}. \emph{Reasoning by Superposition} provided a theoretical foundation for these observations, proving that continuous latent reasoning can encode parallel search states that are inaccessible to discrete CoT \cite{zhu2025reasoning}. Additional studies further explore continuous latent reasoning and hybrid token-latent representations \cite{gozeten2025continuous,su2025token}.

\paragraph{Latent compression and optimization.}
CoLaR accelerates latent reasoning by dynamically compressing chains of thought into fewer latent steps and optimizing reasoning length via reinforcement learning \cite{tan2025think}. In CoLaR, latent representations are explicitly trained to preserve token-level semantics and are optimized for semantic fidelity to compressed textual reasoning trajectories. CODI similarly distills explicit CoT into continuous thoughts via self-distillation \cite{shen2025compressing}, while Token Assorted mixes latent and text tokens for reasoning \cite{su2025token}. While these approaches improve efficiency, they retain autoregressive latent generation and treat latent variables primarily as compact surrogates for language tokens. Consequently, they focus on compressing or distilling existing reasoning traces rather than enabling latent states to reorganize computation or support richer internal reasoning dynamics.

\paragraph{Positioning of our work.}
Our method PaCT builds directly on Coconut’s paradigm of continuous latent reasoning and is guided by the theoretical insights of superposition-based expressivity. In contrast to latent compression approaches such as CoLaR and CODI, we target the architectural source of inefficiency by eliminating recurrent latent unrolling. By synthesizing and refining latent representations in parallel, PaCT preserves the expressive advantages of continuous reasoning while enabling scalable and efficient computation.
